\documentclass[../../main.tex]{subfiles}% 注意这里的文件路径不能用 ./main.tex ,否则用latexmk编译子文件会报错
\graphicspath{{\subfix{./image/}}{\subfix{../../image/}}} % 指定图片引用路径,后续可以直接使用图片文件名
% 如果图片存放在上级目录,则使用 ../../image/ 作为备用路径

% 例如:
% \begin{figure}[H]
% \centering
% \includegraphics[width=0.3\textwidth]{图.png}
% \caption{}
% \label{figure:图}
% \end{figure}
% 注意:上述\label{}一定要放在\caption{}之后,否则引用图片序号会只会显示??.

\begin{document}

\section{多值全纯函数与单值性定理}

\begin{definition}
设幂级数 \(P(z)=\sum\limits_{n=0}^{\infty} a_n(z-a)^n\) 的收敛半径 \(R>0\), \(z=\gamma(t)\) (\(\alpha\leqslant t\leqslant \beta\)) 是以 \(a\) 为起点的平面曲线. 若存在幂级数族 \(P_t(z)=\sum\limits_{n=0}^{\infty} a_n(t)(z-\gamma(t))^n\), 它具有下列性质:
\begin{enumerate}[(1)]
\item 对任意 \(t\in[\alpha,\beta]\), \(P_t\) 的收敛圆盘 \(B_t\) 不退化为一点;
\item 对任意 \(t_0\in[\alpha,\beta]\), 存在 \(\delta>0\), 当 \(t\in(t_0-\delta,t_0+\delta)\cap[\alpha,\beta]\) 时, 总有 \(B_t\cap B_{t_0}\neq\varnothing\), 并且 \(P_t\) 和 \(P_{t_0}\) 在 \(B_t\cap B_{t_0}\) 上恒等;
\item \(P_\alpha=P\),
\end{enumerate}
则称 \(P\) 能沿曲线 \(\gamma\) \textbf{全纯开拓}, 并称 \(P_\beta\) 是 \(P\) 沿曲线 \(\gamma\) 的\textbf{全纯开拓}.
\end{definition}

\begin{examples}
设幂级数 \(P(z)=\sum\limits_{n=0}^{\infty} a_n(z-a)^n\) 的收敛半径 \(R>0\), 若 \(z_0\in \overline{B(a,R)}\) 不是 \(P\) 的奇点, 则 \(P\) 能沿线段 \([a,z_0]\) 全纯开拓; 若 \(w_0\in \partial B(a,R)\) 是 \(P\) 的奇点, 则 \(P\) 不能沿半径 \([a,w_0]\) 全纯开拓.
\end{examples}

\begin{examples}
\begin{align*}
\ln|z| + \mathrm{i}\arg z &= \sum_{n=1}^{\infty} \frac{(-1)^{n-1}}{n} (z-1)^n.
\end{align*}
沿 \(\partial B(0,1)\) 正向的全纯开拓为
\begin{align*}
\ln|z| + \mathrm{i}(\arg z + 2\pi) &= 2\pi\mathrm{i} + \sum_{n=1}^{\infty} \frac{(-1)^{n-1}}{n} (z-1)^n;
\end{align*}
沿 \(\partial B(0,1)\) 负向的全纯开拓为
\begin{align*}
\ln|z| + \mathrm{i}(\arg z - 2\pi) &= - 2\pi\mathrm{i} + \sum_{n=1}^{\infty} \frac{(-1)^{n-1}}{n} (z-1)^n.
\end{align*}
\end{examples}
\begin{remark}
这表明幂级数沿具有相同起点和终点的两条不同曲线的全纯开拓可能不同.
\end{remark}

\begin{proposition}\label{proposition:沿曲线的全纯开拓唯一}
若幂级数能沿某曲线全纯开拓, 则它沿该曲线的全纯开拓是唯一的.
\end{proposition}
\begin{proof}
设 \(P(z)=\sum\limits_{n=0}^{\infty} a_n(z-a)^n\) 的收敛半径 \(R>0\), \(z=\gamma(t)\) (\(\alpha\leqslant t\leqslant \beta\)) 是以 \(a\) 为起点的平面曲线, 幂级数族 \(P_t\) 和 \(Q_t\) 给出了 \(P\) 沿曲线 \(\gamma\) 的两个全纯开拓 \(P_\beta\) 和 \(Q_\beta\), 我们的目标是证明 \(P_\beta=Q_\beta\).

设 \(P_t\) 和 \(Q_t\) 的收敛圆盘分别为 \(B_t\) 和 \(U_t\), \(\Gamma=\{t\in[\alpha,\beta]: P_t=Q_t\}\). 首先, 可证明 \(\Gamma\) 是 \([\alpha,\beta]\) 中的开集. 事实上, 若 \(t_0\in\Gamma\), 则存在 \(\delta>0\), 当 \(t\in(t_0-\delta,t_0+\delta)\cap[\alpha,\beta]\) 时, 成立 \(B_t\cap B_{t_0}\neq\varnothing, P_t=P_{t_0}\Big|_{B_t\cap B_{t_0}}\) 和 \(U_t\cap U_{t_0}\neq\varnothing, Q_t=Q_{t_0}\Big|_{U_t\cap U_{t_0}}\). 因为 \(P_{t_0}=Q_{t_0}, B_{t_0}=U_{t_0}\), 由全纯函数的内部唯一性定理便知 \(P_t=Q_t\), 即 \(t\in\Gamma\).

其次, 可证明 \(\Gamma\) 是 \([\alpha,\beta]\) 中的闭集. 事实上, 若 \(t_0\) 是 \(\Gamma\) 的极限点, 则可取 \(t\in\Gamma\), 使得 \(B_t\cap B_{t_0}\neq\varnothing, P_t=P_{t_0}\Big|_{B_t\cap B_{t_0}}\) 和 \(U_t\cap U_{t_0}\neq\varnothing, Q_t=Q_{t_0}\Big|_{U_t\cap U_{t_0}}\). 因为 \(P_t=Q_t, B_t=U_t\), 由全纯函数的内部唯一性定理便知 \(P_{t_0}=Q_{t_0}\), 即 \(t_0\in\Gamma\).

注意到 \(\Gamma\neq\varnothing\), 便知 \(\Gamma=[\alpha,\beta]\). 因此, \(P_\beta=Q_\beta\).
\end{proof}

\begin{definition}
设 \(D\) 是域, \(f\) 是 \(D\) 上的多值函数, 即对任意 \(z\in D\), \(f(z)\) 是一个非空复数集. 若存在以 \(a\in D\) 为收敛圆心、以 \(R>0\) 为收敛半径的幂级数 \(P_0\), 使得
\begin{enumerate}[(1)]
\item \(P_0\) 能沿 \(D\) 中以 \(a\) 为起点的任意曲线 \(\gamma\) 全纯开拓;
\item 当 \(z\in D\) 时, \(f(z)=\{P(t); P \text{ 是 } P_0 \text{ 沿 } D \text{ 中连接 } a \text{ 和 } z \text{ 的曲线 } \gamma \text{ 的全纯开拓}\}\),
\end{enumerate}
即 \(f\) 这个多值函数是由幂级数 \(P_0\) 在 \(D\) 上经过全纯开拓得到的, 则称 \(f\) 是 \(D\) 上的\textbf{多值全纯函数}.
\end{definition}

\begin{theorem}[单值性定理]\label{theorem:单值性定理}
设 \(f\) 是域 \(D\) 上的多值全纯函数, \(z_0\in D, w_0\in f(z_0)\). 若 \(G\subseteq D\) 是单连通域, 则必能在 \(G\) 上选出 \(f\) 的一个单值全纯分支 \(g\), 使得 \(g(z_0)=w_0\), 并且 \(f\) 可由 \(\sum\limits_{n=0}^{\infty} \frac{g^{(n)}(z_0)}{n!}(z-z_0)^n\) 在 \(D\) 上经过全纯开拓得到.
\end{theorem}
\begin{proof}

设 \(f\) 是由幂级数 \(P_0\) 在 \(D\) 上经过全纯开拓得到的, 故可取 \(D\) 中连接 \(P_0\) 的收敛圆心和 \(z_0\) 的曲线 \(\gamma\), 使得 \(P_0\) 沿 \(\gamma\) 的全纯开拓 \(P\) 满足 \(P(z_0)=w_0\). 注意, \(f\) 也可由幂级数 \(P\) 在 \(D\) 上经过全纯开拓而得到.

不妨设 \(G\) 异于 \(\mathbb{C}\), 否则 \(P\) 的收敛半径为 \(\infty\), 定理显然成立. 由 \hyperref[theorem:Riemann映射定理]{Riemann映射定理}, 存在双全纯映射 \(\varphi\), 使得 \(\varphi(B(0,1))=G, \varphi(0)=z_0\), 故 \(P\circ \varphi(\zeta)\) 在 \(0\) 处全纯. 因此, \(h(\zeta)=\sum\limits_{n=0}^{\infty} \frac{(P\circ \varphi)^{(n)}(0)}{n!} \zeta^n\) 的收敛半径 \(\rho>0\). 我们断言, 对任意 \(\zeta_0 \in B(0,1)\), \(h\) 能沿 \([0,\zeta_0]\) 全纯开拓. 实际上, 对于 \(G\) 中的曲线 \(z=\gamma(t)=\varphi(t\zeta_0)\) (\(0\leqslant t\leqslant 1\)), \(P\) 能沿 \(\gamma\) 全纯开拓. 设幂级数族 \(Q_t\) 给出了 \(P\) 沿 \(\gamma\) 的全纯开拓, 则幂级数族 \(h_t(\zeta)=\sum\limits_{n=0}^{\infty} \frac{(Q_t \circ \varphi)^{(n)}(t\zeta_0)}{n!} (\zeta-t\zeta_0)^n\) 便给出了 \(h\) 沿 \([0,\zeta_0]\) 的全纯开拓. 这表明 \(\zeta_0\) 不是 \(h\) 的奇点, 因此 \(\rho\geqslant 1\). 于是 \(g=h\circ \varphi^{-1}\) 在 \(G\) 上全纯, 并且
\begin{align*}
P(z)=\sum_{n=0}^{\infty} \frac{g^{(n)}(z_0)}{n!} (z-z_0)^n.
\end{align*}
\end{proof}

\begin{theorem}
设 \(f\) 是整函数, 若存在 \(a,b\in\mathbb{C}\), \(a\neq b\), 使得 \(f(\mathbb{C})\subseteq \mathbb{C}\setminus\{a,b\}\), 则 \(f\) 是常值函数.
\end{theorem}
\begin{proof}
不妨设 \(a=0,b=1\), 否则考虑 \(\frac{f(z)-a}{b-a}\). 由 \hyperref[theorem:Riemann映射定理]{Riemann映射定理}、\hyperref[theorem:边界对应定理]{边界对应定理}和\cref{proposition:复变例1112}知, 存在 \(B(0,1)\) 上的全纯函数 \(\varphi\), 使得 \(\varphi(B(0,1))= \mathbb{C}\setminus\{0,1\}\). 再由\hyperref[theorem:单值性定理]{单值性定理} 知, 多值全纯函数 \(\varphi^{-1} \circ f\) 可在 \(\mathbb{C}\) 上选出一个单值全纯分支 \(\psi\). 由 \hyperref[theorem:Liouville(刘维尔)定理]{Liouville定理}, \(\psi\) 是常值函数, 故 \(f\) 也是常值函数. \end{proof}





\end{document}